\chapter{Repaso teórico razonado y preguntas de desarrollo}
\label{ch:repaso-teorico-matmo}
Estas cuestiones son originales y sintetizan razonamientos de los capítulos precedentes, no copian tests de estudiantes. Se recomienda intentar una respuesta propia antes de leer el texto. Las respuestas son guías de estudio, no instrucciones aprobadas ni pruebas oficiales de la ULE.

\section{Materiales y daño}

\paragraph{R01. ¿Por qué una aleación de menor densidad puede generar una pieza más pesada?} Porque la sección necesaria depende de resistencia, rigidez, deformación y vida. Para una barra dimensionada por resistencia y longitud fija, la comparación de masa depende de densidad sobre tensión admisible, no de densidad aislada.

\paragraph{R02. ¿Qué se diferencia entre tensión nominal, tensión local y tenacidad a fractura?} La primera promedia fuerza sobre sección; la segunda incorpora geometría y concentración; la tenacidad representa resistencia a un modo de propagación bajo condiciones definidas. No son propiedades intercambiables.

\paragraph{R03. ¿Qué significa que una grieta central mida 2a?} La longitud total es el doble de a. En la ecuación del factor de intensidad correspondiente se usa la semilongitud a, tras confirmar geometría y factor corrector.

\paragraph{R04. ¿Se puede calcular una vida segura de fatiga solo con la curva S--N?} No: se requieren espectro, tensión media, entallas, ambiente, dispersión, estado superficial, daño y criterios aplicables. Una curva en un régimen es una correlación, no una garantía.

\paragraph{R05. ¿Qué distingue la iniciación de fatiga de la propagación de grieta?} La iniciación describe la formación de una grieta detectable o definida; la propagación calcula su crecimiento a partir de un tamaño inicial. La integral de Paris modela solo un régimen de crecimiento.

\paragraph{R06. ¿Cuál es el error habitual al aplicar Goodman?} Confundir amplitud con máximo o rango. Primero se calculan tensión media, amplitud y razón R; después se adopta el criterio y su campo de validez.

\paragraph{R07. ¿Por qué la fluencia no exige siempre límites de grano?} También operan mecanismos de dislocaciones y difusión dentro de los granos; un monocristal evita ciertos mecanismos intergranulares, no toda deformación dependiente del tiempo.

\paragraph{R08. ¿Qué aporta y qué no aporta Larson--Miller?} Correlaciona tiempo de rotura y temperatura con datos y convenciones determinados. No sustituye ensayos ni justifica extrapolar entre mecanismos de daño diferentes.

\section{Aleaciones, compuestos y ensayos}

\paragraph{R09. ¿De qué familia es el aluminio 8090?} De aleaciones aluminio--litio. Identificarlo como superaleación de níquel es un error de algunos apuntes estudiantiles.

\paragraph{R10. ¿Qué estructuras tienen las fases principales del titanio?} Alfa es hexagonal compacta y beta cúbica centrada en el cuerpo. La aleación concreta y su tratamiento determinan la microestructura final.

\paragraph{R11. ¿Qué diferencia maduración artificial y revenido?} La primera desarrolla precipitados en aleaciones endurecibles; el segundo modifica martensita y otras estructuras de acero tras el temple. Comparten variables tiempo y temperatura, no mecanismo universal.

\paragraph{R12. ¿Pueden tratarse térmicamente los aluminios de fundición?} Sí, existen aleaciones de fundición con tratamiento de solución y envejecimiento, como ciertos estados T6. El defecto de colada y la especificación de producto deben considerarse aparte.

\paragraph{R13. ¿Qué funciones mecánicas realiza la matriz de un compuesto?} Transmite esfuerzo entre refuerzos, contribuye a cargas transversales y cortadura, protege e influye en daño y temperatura. Decir que nunca soporta cargas es falso.

\paragraph{R14. ¿Qué es un laminado GLARE y de qué depende su rigidez?} Combina capas metálicas de aluminio con capas reforzadas con fibra de vidrio. Su respuesta depende de espesores, direcciones y secuencia de apilado.

\paragraph{R15. ¿En qué se diferencian fallo adhesivo, cohesivo y del sustrato?} En el adhesivo se separa la interfaz; en el cohesivo se rompe el adhesivo; en el tercero se rompe el material adherido. El examen exige localizar correctamente la zona.

\paragraph{R16. ¿Se puede elegir radiografía para toda grieta de 5 mm?} No. La profundidad no basta: deben conocerse material, orientación, abertura, geometría, acceso y sensibilidad. Otras técnicas pueden ser más adecuadas en condiciones distintas.

\section{Certificación y regulación aeronáutica}

\paragraph{R17. ¿Qué distingue al reglamento 2018/1139 de 748/2012 y 1321/2014?} El primero establece la base regulatoria y el sistema general; 748/2012 desarrolla aspectos de certificación inicial y diseño/producción; 1321/2014 aborda aeronavegabilidad continuada, mantenimiento y personas u organizaciones implicadas.

\paragraph{R18. ¿Qué acredita un certificado de tipo?} La aprobación del diseño de tipo frente a su base de certificación, en su ámbito. No sustituye el certificado de aeronavegabilidad de una aeronave concreta ni un CRS tras una tarea.

\paragraph{R19. ¿Qué son un STC y una modificación de diseño?} Un certificado de tipo suplementario ampara ciertas modificaciones mayores del diseño de tipo conforme a la vía aplicable. Su existencia no exime comprobar efectos en instrucciones de mantenimiento y configuración.

\paragraph{R20. ¿Qué debe analizarse antes de aplicar Parte ML?} Categoría, masa, complejidad, operación y excepciones del artículo de ámbito. No basta con conocer un solo dato como la masa de la aeronave.

\paragraph{R21. ¿Qué es una CAMO?} Una organización de gestión de aeronavegabilidad continuada dentro de privilegios aprobados. Gestionar aeronavegabilidad no equivale a ejecutar y certificar toda tarea de mantenimiento.

\paragraph{R22. ¿Qué diferencia Parte 145, 66 y 147?} Parte 145 se ocupa de organizaciones de mantenimiento; Parte 66, licencias del personal; Parte 147, organizaciones de formación. Los privilegios de cada una no son automáticos para las demás.

\paragraph{R23. ¿Cuándo puede emitirse un ARC si la revisión revela discrepancias?} El proceso debe satisfacer los requisitos pertinentes y resolver las discrepancias que impidan la emisión. Un certificado anterior o una firma aislada no sustituyen la revisión.

\paragraph{R24. ¿En qué se distinguen el adelanto de una revisión y la extensión de un ARC?} Son procedimientos distintos y con límites condicionados. No debe confundirse el margen admitido en una revisión con el de una extensión ni con tolerancias de tareas de mantenimiento.

\section{Gestión del mantenimiento y casos prácticos}

\paragraph{R25. ¿Qué es un AMP y por qué no es solo un MPD?} El AMP es el programa de mantenimiento aplicable a una aeronave y configuración en su régimen; el MPD puede ser fuente de tareas del fabricante. Documento fuente y programa individual no son equivalentes.

\paragraph{R26. ¿Qué significa FH, FC y límite calendario?} Horas de vuelo, ciclos de vuelo y tiempo de calendario son magnitudes diferentes. Una tarea con varios límites puede vencer por la primera condición alcanzada.

\paragraph{R27. ¿Qué requiere registrar la realización de una tarea?} Identidad de tarea, datos y revisión aplicables, aeronave/componente, fecha, acumulados, resultados, personal y certificación que corresponda. El cómputo del plazo no demuestra ejecución.

\paragraph{R28. ¿Por qué una AD no se sustituye por un SB más reciente?} Porque la AD puede incorporar un documento y revisión concretos, con condiciones de aplicabilidad y método. Una alternativa exige comprobar las opciones admitidas o la aprobación pertinente.

\paragraph{R29. ¿Para qué sirven AMM, CMM, IPC y SRM?} AMM aporta instrucciones de mantenimiento de aeronave, CMM de componente, IPC identificación de piezas y SRM reparaciones estructurales en su ámbito. Hay que comprobar revisión y aplicabilidad.

\paragraph{R30. ¿Qué información demuestra que una LLP sigue dentro de su vida?} Identidad, ciclos u otra magnitud definida, historial trazable, límite aplicable y estado del componente. La ausencia de datos no autoriza reiniciar el contador a cero.

\paragraph{R31. ¿Qué significa una tendencia de EGT deteriorada?} Es un indicador que exige normalizar condiciones y correlacionar otros parámetros. No identifica automáticamente el componente dañado ni autoriza extender vida.

\paragraph{R32. ¿Por qué EMAR no equivale a regulación militar europea directamente obligatoria?} Es una base de armonización promovida en cooperación de defensa; la obligación concreta exige examinar la normativa nacional, autoridad y ámbito. En España se estudia el Reglamento de Aeronavegabilidad de la Defensa.

\section{Comprobación final antes del examen}
Una explicación sólida distingue definición, mecanismo, modelo, hipótesis y limitaciones. En problemas, se comprueban unidades, casos límite y plausibilidad. En normativa, se identifica primero el ámbito y la versión del texto y se razona solo con hechos conocidos. El banco de test y los simulacros de este proyecto se encuentran en el cuaderno complementario; la selección concreta de preguntas del examen queda fuera del conocimiento del manual.
